% Generated draft: requires mathematical and source review.
\documentclass[11pt]{article}
\usepackage[margin=25mm]{geometry}
\usepackage{fontspec}
\setmainfont{Latin Modern Roman}
\usepackage{amsmath,amssymb,mathtools,unicode-math}
\setmathfont{Latin Modern Math}
\usepackage{enumitem,needspace,xurl,hyperref}
\setlength{\parindent}{0pt}
\setlength{\parskip}{5pt}
\setlength{\emergencystretch}{3em}
\newcommand{\N}{\mathbb{N}}
\newcommand{\Z}{\mathbb{Z}}
\newcommand{\Q}{\mathbb{Q}}
\newcommand{\R}{\mathbb{R}}
\newcommand{\C}{\mathbb{C}}
\newcommand{\F}{\mathbb{F}}
\newcommand{\A}{\mathbb{A}}
\newcommand{\PP}{\mathbb{P}}
\newcommand{\E}{\mathbb{E}}
\newcommand{\eps}{\varepsilon}
\newcommand{\dd}{\,\mathrm{d}}
\begin{document}
\section*{UnsolvedMath ID: 11000163}
\textbf{Source problem code:} catalogue entry
\subsection{Definitions and mathematical statement}
For a Heegaard splitting \(M=V_1\cup_\Sigma V_2\) in the standard closed oriented setting, let \(\mathcal D(V_i)\) be the set of essential curves on \(\Sigma\) bounding disks in \(V_i\). For genus at least \(2\), its Hempel distance is
\[
 d(\Sigma)=\min\{d_{\mathcal C(\Sigma)}(\alpha,\beta):
       \alpha\in\mathcal D(V_1),\ \beta\in\mathcal D(V_2)\},
\]
where the curve graph has essential-curve vertices and disjointness edges. Low-genus and boundary conventions require separate specification.

The source reports Schleimer [42] proving that each fixed \(3\)-manifold \(M\) has a finite uniform bound \(B(M)\) on distances of all its Heegaard splittings. The task is to study this bound and develop an algorithm for computing it. With a finite triangulation as input, a natural precise algorithmic goal is to output an integer \(B(M)\) such that
\[
 \forall\Sigma\text{ a Heegaard splitting of }M,\qquad d(\Sigma)\leq B(M).
\]
The source does not say whether it requires any computable bound, the particular theorem's bound, or the optimal bound \(\sup_\Sigma d(\Sigma)\). It imposes no fixed Heegaard genus in the printed statement.

The following background defines the handlebody subgroup \(H_g\) as surface mapping classes extending over a genus-\(g\) handlebody. It credits Wajnryb [51] with a finite presentation, reports simplification only in genus two at that time, and describes the induced subgroup of \(\operatorname{Sp}(2g,\Z)\) on surface homology as transparent while the group's structure is little understood. The extracted \(H_1(2g,\Z)\) notation is not defined; usual surface homology is \(\Z^{2g}\). A further handlebody problem is announced but not printed, so none is reconstructed.
\subsection{Short English statement}
Make the reported uniform Heegaard-distance bound for a fixed \(3\)-manifold algorithmically computable, clarifying whether an arbitrary upper bound or the optimal bound is required.
\subsection{Sources}
\begin{itemize}
\item J. Birman, chapter in \emph{Problems on Mapping Class Groups and Related Topics}, edited by Farb, Problem 2.6 and following handlebody discussion, supplied PDF page 156. \url{https://www.math.uchicago.edu/~farb/papers/mcgbook.pdf}. Recorded access 2026-07-29; extraction pdf-font-structure.
\item Schleimer [42], fixed-manifold distance bound; Wajnryb [51], handlebody-subgroup presentation. Full bibliography and theorem hypotheses are absent. The summary mentions related Hempel work without an exact reference.
\item Local UnsolvedMath registry, \texttt{data/unsolvedmath/problems.json}, permanent ID \(11000163\), exact source code \texttt{catalogue entry}; supplied OPEN-TRIAGE summary distinguishes the known bound from the follow-up computation problem.
\end{itemize}
\end{document}
